\appendix
\section{Post-Project Extension: Self-Supervised Vision Transformer Baselines}
\label{app:dino}

\textit{Explanatory Note: The architectures documented in this appendix were evaluated in an exploratory capacity following the formal conclusion of the primary project evaluation period. Crucially, these exploratory runs were conducted iteratively and carried the identical test-informed evaluation caveat as earlier prototype iterations, prior to the definitive freezing of our strict zero-test-leakage protocol. They are documented here to provide technical completeness, situate our empirical findings relative to the published literature, and outline strategic directions for future development; they did not form part of the benchmarked production scope detailed in the core business report.}

\subsection{Architectural Formulation and Literature Context}

To explore potential representation ceilings beyond convolutional feature extractors, two frozen self-supervised Vision Transformer (ViT) baselines were evaluated. This investigation is directly anchored in recent breakthroughs reported in the industrial anomaly detection literature, most prominently the \textit{Dinomaly} benchmark (Bae et al., 2024). The \textit{Dinomaly} architecture demonstrated that large foundation models pre-trained via discriminative self-distillation can achieve an upper-bound image-level AUROC of 99.6\% and a pixel-level AUROC of 98.2\% on the MVTec AD benchmark when leveraging a frozen DINOv2 \cite{oquab2023dinov2} backbone coupled with task-specific projection heads.

\subsubsection{Self-Supervised Baselines (ViT-S/14)}
We extracted patch-level token representations from the final encoder block of two off-the-shelf Vision Transformer models:
\begin{itemize}[leftmargin=1.5em]
    \item \textbf{DINOv2 (ViT-S/14, frozen):} Patch-level cosine nearest-neighbour matching against an empirical memory bank of frozen self-supervised ViT patch tokens, pre-trained via discriminative self-distillation on the curated LVD-142M dataset.
    \item \textbf{DINOv3 (ViT-S/14, frozen):} An identical pipeline substituting the experimental v3 weights, which extend the self-distillation objective using expanded training sets and refined multi-crop augmentations.
\end{itemize}

Both architectures adhere to the identical non-parametric memory-bank paradigm established by PatchCore -- zero client-side backpropagation and rapid single-pass onboarding -- substituting the convolutional ResNet backbone with a multi-head self-attention transformer encoder, without introducing the complex fine-tuning heads utilised in \textit{Dinomaly}.

\subsection{Indicative Empirical Results and Engineering Trade-Offs}

\begin{table}[b!]
\centering
\caption{DINOv2 and DINOv3 (ViT-S) macro-averaged exploratory benchmark performance. Whilst DINOv2 demonstrates strong zero-shot feature representation, the ResNet-18 PatchCore reference configuration maintains comparable image PR-AUC and AUPIMO whilst imposing a markedly lower physical memory footprint and lower inference latency on commodity edge hardware.}
\label{tab:dino_results}
\scriptsize
\begin{tabular}{@{}lcccc@{}}
\toprule
\textbf{Architecture} & \textbf{Image $F_1$ $\uparrow$} & \textbf{Image PR-AUC $\uparrow$} & \textbf{AUPIMO $\uparrow$} & \textbf{Pixel AUROC $\uparrow$} \\
\midrule
DINOv2 (ViT-S/14, frozen) & 0.941 & 0.989 & 0.602 & 0.969 \\
DINOv3 (ViT-S/14, frozen) & 0.957 & 0.993 & 0.680 & 0.967 \\
\midrule
\textit{PatchCore ResNet-18 (reference)} & \textit{0.929} & \textit{0.988} & \textit{0.603} & \textit{0.968} \\
\bottomrule
\end{tabular}
\end{table}

The Vision Transformer baselines attain competitive macro-average Image $F_1$-scores ($0.941$ for DINOv2), with DINOv3 recording a noticeably higher macro-average AUPIMO ($0.680$ versus $0.603$). Relative to the published \textit{Dinomaly} ceiling ($98.2\%$ pixel AUROC), our direct patch-token extraction achieves $96.9\%$ without requiring auxiliary decoders or heavy task adaptation, placing our baseline just $1.3$ points away from the state-of-the-art pixel ceiling.

Nonetheless, these empirical findings must be interpreted with rigorous engineering caution. First, as disclosed above, these exploratory DINO experiments were evaluated iteratively and carried the test-informed caveat; they were not subjected to the pristine, single-pass protocol lock enforced on the primary production benchmark. Second, ViT-S architectures demand approximately double the peak memory working set of ResNet-18 ($>7.5\,\text{GB}$ vs $3.8\,\text{GB}$). Furthermore, unlike convolutions which benefit heavily from spatial locality and AVX-512 vectorisation on standard edge CPUs, the global self-attention mechanism creates severe compute bottlenecks during feature extraction, drastically inflating the inference latency.

Consequently, whilst Vision Transformers represent a highly promising medium-term research trajectory (particularly when adapted via dedicated decoders as in \textit{Dinomaly}), the frozen ResNet-18 PatchCore engine remains our unequivocal recommendation for immediate plant-floor deployment. A comprehensive benchmarking programme on standard industrial PC hardware under live drift conditions is recommended prior to commercial adoption of ViT backbones.
